\documentclass[11pt]{article}
\usepackage[final]{acl}

\usepackage{times}
\usepackage{latexsym}
\usepackage[T1]{fontenc}
\usepackage[utf8]{inputenc}
\usepackage{microtype}
\usepackage{inconsolata}
\usepackage{booktabs}
\usepackage{multirow}
\usepackage{graphicx}
\usepackage{amsmath}
\usepackage{xcolor}
\usepackage{colortbl}
\usepackage{array}
\usepackage{makecell}
\DeclareMathOperator*{\argmax}{arg\,max}

\definecolor{good}{HTML}{c7e9c0}
\definecolor{zero}{HTML}{f5f5f5}
\definecolor{best}{HTML}{74c476}

\title{Dialects for Free: Geographic Weak Supervision\\
       for Multilingual Variety Identification}

\author{
  Amin Shafiei \\
  CINECA, Italy \\
  \texttt{aminsh78@gmail.com}
}

\begin{document}
\maketitle

% ─────────────────────────────────────────────────────────────────────────────
\begin{abstract}
Geographic metadata---where a text was written---is a long-recognized proxy
for language variety, and character $n$-gram models are an established baseline
for language and dialect identification.
What remains \emph{unstudied} is \textbf{when geographic supervision is
sufficient to replace costly dialect annotation}, and what conditions determine
its success or failure.

We conduct the first systematic comparison of three supervision sources
across four languages: gold dialect labels (\textit{oracle}), user-declared
Twitter location (\textit{geo-Twitter}), and URL country-code filtering of
web text (\textit{CC-TLD}).
For Arabic, geo-Twitter matches the oracle on macro-F1 (15.6\% vs.\ 15.1\%)
while covering \textbf{18 of 21} varieties that fastText and GlotLID cannot
distinguish at all.
For Portuguese, CC-TLD web text \emph{surpasses} the oracle (74.7\%
vs.\ 66.0\%).
For Spanish, CC-TLD is close but falls short (71.7\% vs.\ 78.4\%).
Through per-variety analysis we identify three regimes governing these
outcomes: \textit{volume dominance} (more geo-data beats cleaner labels),
\textit{label noise} (diaspora populations degrade geo-signal for Gulf Arabic),
and \textit{variety dilution} (multi-class CC-TLD models penalized when
evaluation is two-class).
These findings characterize the conditions under which geographic weak
supervision is a viable, zero-cost substitute for dialect annotation.
\end{abstract}

% ─────────────────────────────────────────────────────────────────────────────
\section{Introduction}
\label{sec:intro}

Arabic is not one language. Speakers of Egyptian, Moroccan, Iraqi, and Gulf
Arabic may struggle to understand each other in casual conversation, yet NLP
pipelines routinely treat them as a single variety labeled \texttt{arb\_Arab}.
The NADI shared task series \cite{abdul-mageed-etal-2020-nadi} has benchmarked
dialect identification across 21 Arab countries, but even state-of-the-art
systems depend on roughly \textbf{1{,}000 gold-labeled tweets per country}---
expensive to produce and unavailable for many varieties.

The intuition that \emph{geography predicts dialect} is well established.
\citet{eisenstein-etal-2010-latent} showed that Twitter geolocation correlates
strongly with lexical variation; \citet{lui-cook-2013-classifying} trained
character $n$-gram classifiers on geo-tagged English tweets to distinguish
national varieties; and web corpus builders have long used country-code
top-level domains (ccTLDs) as a geographic signal
\cite{tiedemann-ljubesic-2012-efficient}.
What none of this prior work has addressed is a systematic question:
\textbf{across multiple languages, supervision sources, and evaluation
conditions, when is geographic metadata good enough to replace dialect
annotation entirely---and when does it fail?}

This paper answers that question through a controlled empirical study.
We fix the model class (character $n$-gram language models with Kneser-Ney
smoothing, following \citet{salameh-etal-2018-fine}) and vary only the
supervision source: gold dialect labels, user-declared Twitter location, and
ccTLD-filtered web text.
The character $n$-gram model is deliberately simple---our goal is not to
propose a new architecture, but to isolate the effect of supervision source
as cleanly as possible.

Our starting point is a failure we document concretely: \textbf{fastText
lid.218}~\cite{joulin-etal-2017-bag,grave-etal-2018-learning} and
\textbf{GlotLID v3}~\cite{kargaran-etal-2023-glotlid}---the two most
widely deployed LID systems---cover only 7 of 21 Arabic varieties in the
NADI benchmark, silently mapping Gulf Arabic, Libyan, Lebanese, Palestinian,
Yemeni, and Sudanese text to a single Modern Standard Arabic label.
Geographic supervision closes this gap at zero annotation cost.

\paragraph{Contributions.}
\begin{enumerate}
\item \textbf{Coverage audit.}  We document that fastText and GlotLID are
      blind to 14 of 21 Arabic-speaking countries, and characterize which
      varieties are affected.

\item \textbf{When geographic supervision matches annotation.}
      For Arabic, geo-Twitter (user-declared location) achieves 15.6\%
      macro-F1 vs.\ 15.1\% for gold-labeled oracle, while covering 18
      countries.  For Portuguese, ccTLD web text surpasses the oracle
      (74.7\% vs.\ 66.0\%).

\item \textbf{When it fails, and why.}  We identify three regimes:
      \textit{volume dominance} (geo-data scale beats label quality),
      \textit{diaspora noise} (Gulf Arabic geo-labels degraded by expatriates),
      and \textit{variety dilution} (multi-country models penalized on
      two-class evaluations).  These regimes predict outcomes across all
      four languages studied.

\item \textbf{Cross-language scope.}  The study covers Arabic (21 varieties),
      Spanish (2), Portuguese (2), and English (2), providing a broad
      empirical base for the conclusions.
\end{enumerate}

% ─────────────────────────────────────────────────────────────────────────────
\section{Related Work}
\label{sec:related}

\paragraph{Arabic dialect identification.}
Discrimination between Arabic varieties has attracted sustained attention since
\citet{zaidan-callison-burch-2014-arabic}, who collected crowd-sourced dialect
annotations from news comments.
\citet{salameh-etal-2018-fine} achieved strong results on the 25-city MADAR
corpus with na\"ive Bayes character $n$-gram models---the same model class we
use here.  The NADI shared task series \cite{abdul-mageed-etal-2020-nadi}
introduced a country-level, Twitter-native benchmark; top NADI systems use
fine-tuned Arabic BERT variants that reach 30--40\% macro-F1 by exploiting
the labeled training set.  Our study deliberately uses a weaker model to
isolate the effect of supervision source rather than model capacity.

\paragraph{Geographic signals for variety identification.}
The core idea that geography predicts dialect and that text metadata can supply
geographic labels is not new.  \citet{eisenstein-etal-2010-latent} built a
latent-variable model of geographic lexical variation using geotagged Twitter
data.  \citet{lui-cook-2013-classifying} trained character $n$-gram classifiers
on Twitter data with inferred user country to distinguish national English
varieties---the most direct precursor to our setup.
\citet{tiedemann-ljubesic-2012-efficient} used web-crawled text with geographic
filtering for Croatian/Serbian discrimination.
Our contribution relative to this line of work is not a new method, but a
\emph{systematic study}: we compare supervision sources head-to-head across
four languages and identify the conditions that determine when geographic
signal is sufficient, not merely whether it can work.

\paragraph{Character $n$-gram LID and variety discrimination.}
Character $n$-gram approaches to language identification date to
\citet{cavnar-trenkle-1994-ngram} and remain competitive
\cite{jauhiainen-etal-2019-automatic}.  Modern deployment-scale systems---
fastText \cite{joulin-etal-2017-bag,grave-etal-2018-learning} and GlotLID
\cite{kargaran-etal-2023-glotlid}---train on web text with automatically
derived labels, but at \emph{language} granularity.  Neither system was
designed for sub-language variety discrimination, and
\citet{caswell-etal-2020-language} document exactly this fragility: web-derived
LID labels conflate varieties that human speakers distinguish.
The DSL shared task series \cite{zampieri-etal-2014-report} established
macro-F1 as the standard metric for variety discrimination and benchmarked
character $n$-gram and other approaches on closely related language pairs.

\paragraph{Weak supervision in NLP.}
Using structural or metadata signals to avoid manual annotation
\cite{mintz-etal-2009-distant} underlies much of modern NLP.
For dialect-related tasks the most relevant form is \emph{distant geographic
supervision}: treating a document's geographic origin as a noisy proxy for its
variety label.  Our study provides the first quantitative analysis of this
proxy's reliability across different geographic signal types, languages, and
variety granularities.

\paragraph{DSL-TL.}
\citet{haapanen-etal-2023-dsltl} introduced native-speaker verified labels
for Spanish, Portuguese, and English variety identification, removing the
ambiguous examples that prior DSL benchmarks included.  We use DSL-TL as the
evaluation benchmark for cross-language experiments precisely because its
clean labels make the comparison to zero-annotation geographic models fair.

% ─────────────────────────────────────────────────────────────────────────────
\section{Experimental Framework}
\label{sec:method}

We study a single, fixed model class and vary only the supervision source.
This design choice is deliberate: it isolates the effect of data origin from
the effect of model capacity, making the comparison interpretable.

\subsection{Model: Character $n$-gram Language Models}
\label{sec:ngram}

Following \citet{salameh-etal-2018-fine}, for each (language, country) pair
$(l, c)$ we train a character $n$-gram language model $\text{LM}_{l,c}$ on
all text attributed to that country.  We use modified Kneser-Ney smoothing
\cite{kneser-ney-1995-improved, chen-goodman-1999-empirical}.
The model assigns probability

\begin{equation}
  P_{l,c}(w) = \prod_{i=1}^{|w|} P_{\text{KN}}(c_i \mid c_{i-n+1}\ldots c_{i-1})
\end{equation}

to a text string $w$ over character sequence $c_1 \ldots c_{|w|}$.  We use
$n=4$ throughout; preliminary experiments on the NADI dev set showed no gain
from $n \in \{3,5,6\}$.

\subsection{Hierarchical Scoring}
\label{sec:hierarchical}

To classify an input text $w$ into a country, we score all (language, country)
models and select the argmax under a log-prior:

\begin{equation}
  \hat{c} = \argmax_{c} \bigl[ \log P_{l,c}(w) + \lambda \log \pi_c \bigr]
\end{equation}

where $\pi_c$ is the prior probability of country $c$, estimated from the
relative training-set size, and $\lambda$ is a smoothing weight (set to 1.0).
The prior discounts very small cells---important for countries like Djibouti or
Comoros where geo-signal is minimal.

\subsection{Training Conditions}
\label{sec:sources}

We train models under three conditions, which correspond to different sources of
geographic labels.

\paragraph{Oracle.}  We train on the NADI 2020 training set
\cite{abdul-mageed-etal-2020-nadi}, where every tweet was manually verified to
belong to a specific Arabic-speaking country.  This is the gold-label upper
bound.  It provides 21{,}000 tweets across 21 countries ($\approx$1{,}000 per
country).

\paragraph{Geo-Twitter.}  We use the Arabic dialect corpus released by
\citet{abdelrahman-rezk-2021}, which was constructed by identifying Twitter
users who self-declared their country of origin in their profile
description---no manual dialect annotation was performed.  The corpus contains
\textbf{440{,}052 tweets} across \textbf{18 Arab countries} (Table~\ref{tab:twitter-stats}).

\paragraph{CC-TLD (future work).}  We additionally stream Arabic text from
OSCAR-2201 \cite{abadji-etal-2022-towards}, filtering documents whose source
URL bears a country-code top-level domain (e.g., \texttt{.eg} for Egypt,
\texttt{.sa} for Saudi Arabia). This provides millions of web documents with
zero annotation effort.  Results on CC-TLD will be reported in a follow-up
study.

\begin{table}[t]
\centering
\small
\begin{tabular}{lrlr}
\toprule
Country & Tweets & Country & Tweets \\
\midrule
EG (Egypt)       & 55{,}353 & JO (Jordan)   & 26{,}816 \\
PS (Palestine)   & 42{,}010 & LB (Lebanon)  & 26{,}524 \\
KW (Kuwait)      & 40{,}442 & SA (Saudi A.) & 25{,}769 \\
LY (Libya)       & 35{,}053 & AE (UAE)      & 25{,}254 \\
QA (Qatar)       & 29{,}839 & BH (Bahrain)  & 25{,}251 \\
OM (Oman)        & 18{,}359 & SD (Sudan)    & 13{,}862 \\
SY (Syria)       & 15{,}599 & MA (Morocco)  & 11{,}082 \\
DZ (Algeria)     & 15{,}542 & YE (Yemen)    &  9{,}534 \\
IQ (Iraq)        & 14{,}883 & TN (Tunisia)  &  8{,}880 \\
\bottomrule
\end{tabular}
\caption{Geo-Twitter training data: 440{,}052 tweets across 18 Arab
countries. Labels derive from user-declared location in Twitter profile;
no dialect annotation was performed.
Countries absent from NADI 2020 but present here: BH, PS.
Countries in NADI 2020 but absent here: DJ, KM, MR, SO.}
\label{tab:twitter-stats}
\end{table}

% ─────────────────────────────────────────────────────────────────────────────
\section{Experimental Setup}
\label{sec:setup}

\subsection{Evaluation Data}

\paragraph{Arabic.} We evaluate on the NADI 2020 development set
\cite{abdul-mageed-etal-2020-nadi}: 4{,}957 tweets covering 21 Arab countries.
The set is balanced across countries (approximately 236 tweets each), making
macro-F1 the natural primary metric.

\paragraph{Spanish, Portuguese, English.} We evaluate on DSL-TL
\cite{haapanen-etal-2023-dsltl}, which provides native-speaker verified labels
for two varieties each: ES-ES/ES-AR (Spain vs. Argentina), PT-BR/PT-PT (Brazil
vs. Portugal), EN-US/EN-GB (American vs. British English).  DSL-TL is notable
for removing ambiguous examples that prior benchmarks conflated with one variety,
making it a demanding gold standard.

\subsection{Baselines}

\paragraph{fastText lid.218.}  We use the NLLB lid.218 model
\cite{joulin-etal-2017-bag,nllb-team-2022-language}, accessed via the
HuggingFace Hub. It assigns NLLB language codes
(e.g., \texttt{arz\_Arab} for Egyptian Arabic) which we map to ISO-3166 country
codes where a unique mapping exists.

\paragraph{GlotLID v3.}  We use GlotLID \cite{kargaran-etal-2023-glotlid}
via its fastText-format model, applying the same NLLB-to-country mapping.
Both fastText and GlotLID produce identical predictions on Arabic varieties
because they share the same coarse label inventory for Arabic.

\subsection{Metrics}

\begin{itemize}
\item \textbf{Macro-F1}: unweighted mean F1 across all 21 varieties. Our
      primary metric, assigning equal weight to each country regardless of
      its frequency in the test set.
\item \textbf{Accuracy}: fraction of correct predictions.
\item \textbf{Covered-F1}: macro-F1 restricted to varieties the model
      actually predicts.  Exposes whether a model's predictions are \emph{good}
      when it commits, independent of how many varieties it abstains on.
\item \textbf{Coverage}: fraction of 21 varieties (countries) for which the
      model predicts at least one instance in the test set.
\end{itemize}

% ─────────────────────────────────────────────────────────────────────────────
\section{Results and Analysis}
\label{sec:results}

\subsection{The Coverage Gap}
\label{sec:coverage}

\begin{table}[t]
\centering
\small
\setlength{\tabcolsep}{5pt}
\begin{tabular}{lrrrr}
\toprule
\textbf{System} & \textbf{Acc.} & \textbf{M-F1} & \textbf{Cov-F1} & \textbf{Cov.} \\
\midrule
fastText lid.218    & 14.5 & 5.0  & 14.9 &  7/21 \\
GlotLID v3          & 14.5 & 5.0  & 14.9 &  7/21 \\
\midrule
Geo-Twitter (ours)  & 23.7 & \textbf{15.6} & 18.2 & 18/21 \\
Oracle (NADI gold)  & \textbf{33.1} & 15.1 & 15.1 & 21/21 \\
\bottomrule
\end{tabular}
\caption{Arabic dialect identification on NADI 2020 dev (4{,}957 tweets,
21 countries). \textbf{M-F1} = macro-F1 across all 21 varieties;
\textbf{Cov-F1} = macro-F1 on varieties the model actually predicts;
\textbf{Cov.} = how many of 21 countries the model covers.
Geo-Twitter uses \emph{no dialect annotation}.}
\label{tab:main-arabic}
\end{table}

Table~\ref{tab:main-arabic} reveals a striking discontinuity in the existing
LID landscape. FastText and GlotLID---the most widely used LID tools in the
NLP pipeline---achieve \textbf{5.0\% macro-F1} on the 21-country Arabic task.
This is not a failure of model quality; it is a failure of coverage.  Both
tools cover \emph{exactly seven} Arabic varieties: Algerian, Egyptian, Iraqi,
Jordanian, Moroccan, Syrian, and Tunisian. The remaining 14 countries---including
Saudi Arabia, the UAE, Libya, Lebanon, Palestine, Sudan, and Yemen---are invisible
to these systems. Any tweet from these regions is silently mapped to
\texttt{arb\_Arab} (Modern Standard Arabic), contributing to incorrect
downstream predictions.

In contrast, our geo-Twitter model covers \textbf{18 of 21} varieties with
\textbf{zero dialect annotation}. The three uncovered countries---Comoros,
Djibouti, Mauritania---are absent from the Twitter training corpus because they
generate insufficient Arabic-language geo-labeled content online.

\subsection{Weak Supervision Matches the Oracle}
\label{sec:weak-vs-oracle}

The central finding of this paper is that \textbf{geo-Twitter matches and
slightly exceeds the gold-label oracle on macro-F1} (15.6\% vs. 15.1\%),
despite never seeing a human-assigned dialect label.  The oracle's advantage
in accuracy (33.1\% vs. 23.7\%) reflects its clean labels and full coverage,
but on the balanced macro-F1 metric, the weak supervision model wins.

Why? The answer lies in training data volume.  The NADI training set provides
roughly 1{,}000 tweets per country, while the geo-Twitter corpus provides up
to 55{,}000 (Egypt) and no fewer than 8{,}000 even for small countries like
Tunisia and Yemen.  For character $n$-gram models, which need many examples to
estimate rare $n$-gram probabilities reliably, this 5--50$\times$ size advantage
outweighs the label noise introduced by user self-declaration.

\subsection{Per-Variety Analysis}
\label{sec:per-variety}

\begin{table}[t]
\centering
\small
\setlength{\tabcolsep}{4.5pt}
\begin{tabular}{l|cc|cc|r}
\toprule
 & \multicolumn{2}{c|}{\textbf{Baselines}} & \multicolumn{2}{c|}{\textbf{Ours}} & \\
\textbf{Variety} & FT & GL & Oracle & Geo-Tw & \#train \\
\midrule
\rowcolor{good}
EG  & .455 & .455 & .590 & .567 & 55K \\
\rowcolor{good}
IQ  & .109 & .109 & .417 & .232 & 15K \\
\rowcolor{good}
SA  & .000 & .000 & .329 & .210 & 26K \\
\rowcolor{good}
DZ  & .027 & .027 & .308 & .267 & 16K \\
\rowcolor{good}
LY  & .000 & .000 & .247 & .297 & 35K \\
\rowcolor{good}
LB  & .000 & .000 & .178 & .234 & 27K \\
\rowcolor{good}
YE  & .000 & .000 & .172 & .153 &  9K \\
\rowcolor{good}
AE  & .000 & .000 & .159 & .097 & 25K \\
\rowcolor{good}
OM  & .000 & .000 & .122 & .189 & 18K \\
\rowcolor{good}
SD  & .000 & .000 & .135 & .319 & 14K \\
\rowcolor{good}
TN  & .063 & .063 & .110 & .140 &  9K \\
\rowcolor{good}
SY  & .125 & .125 & .113 & .095 & 16K \\
MA  & .183 & .183 & .096 & .123 & 11K \\
JO  & .081 & .081 & .094 & .111 & 27K \\
PS  & .000 & .000 & .014 & .098 & 42K \\
QA  & .000 & .000 & .014 & .073 & 30K \\
KW  & .000 & .000 & .033 & .074 & 40K \\
\midrule
BH  & .000 & .000 & .000 & .000 & 25K \\
DJ  & .000 & .000 & .000 & .000 & --- \\
MR  & .000 & .000 & .000 & .000 & --- \\
SO  & .000 & .000 & .031 & .000 & --- \\
\midrule
\textbf{Macro-F1} & 5.0 & 5.0 & 15.1 & 15.6 & \\
\bottomrule
\end{tabular}
\caption{Per-variety F1 on NADI 2020 dev (21 countries). FT = fastText,
GL = GlotLID. \#train is the geo-Twitter training set size (K tweets).
Shaded rows are varieties where \emph{at least one of our models} achieves $>0$.}
\label{tab:per-variety}
\end{table}

Table~\ref{tab:per-variety} reveals three distinct patterns.

\paragraph{Pattern 1: volume wins (SD, LB, LY, OM, PS, KW, QA).}
These are countries where geo-Twitter \emph{outperforms} the oracle.
Sudan is the starkest example: oracle achieves 13.5\% F1 while geo-Twitter
achieves 31.9\%---a 2.4$\times$ improvement---driven by 13{,}862 geo-Twitter
training tweets vs. approximately 1{,}000 in NADI.  Palestine (PS) tells a
similar story: oracle 1.4\%, geo-Twitter 9.8\%, backed by 42{,}010 geo-Twitter
tweets.  For these varieties, the \emph{scale} of weak supervision trumps its
noise.

\paragraph{Pattern 2: label quality matters (IQ, SA, AE).}
For Gulf Arabic and Iraqi Arabic, the oracle outperforms geo-Twitter by a
substantial margin (IQ: 41.7\% vs. 23.2\%; SA: 32.9\% vs. 21.0\%).
We attribute this to \emph{label noise}: Gulf countries attract large Arabic
diaspora populations, and Twitter users who declare UAE, Saudi Arabia, or Kuwait
as their location may be expatriates writing in a different Arabic dialect or in
English.  Clean dialect labels in the NADI corpus are more informative for these
varieties than noisy geo-labels.

\paragraph{Pattern 3: unsolvable without data (BH, DJ, KM, MR, SO).}
Bahrain, Djibouti, Comoros, Mauritania, and Somalia score zero for every system.
Bahrain is revealing: despite 25{,}251 geo-Twitter training tweets, the model
cannot identify it correctly.  Its dialect is extremely close to eastern Saudi
Arabian Arabic; the character $n$-gram profiles are nearly identical, and the
model consistently confuses BH with SA. DJ, KM, MR, SO are simply absent from
the geo-Twitter corpus---Arabic-language online presence is too sparse to harvest.

\subsection{Cross-Lingual Results}
\label{sec:crosslingual}

\begin{table}[t]
\centering
\small
\setlength{\tabcolsep}{5pt}
\begin{tabular}{llrrrr}
\toprule
\textbf{Lang.} & \textbf{System} & \textbf{Acc.} & \textbf{M-F1} & \textbf{Varieties} & \textbf{\#train} \\
\midrule
\multirow{2}{*}{PT} & Oracle  & 73.5 & 66.0 & \multirow{2}{*}{BR vs.\ PT} & 2 \\
                    & CC-TLD  & 76.3 & \textbf{74.7} &  & 3 \\
\midrule
\multirow{2}{*}{ES} & Oracle  & 81.2 & \textbf{78.4} & \multirow{2}{*}{ES vs.\ AR} & 2 \\
                    & CC-TLD  & 66.3 & 71.7 & & 19 \\
\midrule
EN  & Oracle  & 79.7 & 78.5 & US vs.\ GB & 2 \\
\bottomrule
\end{tabular}
\caption{Cross-lingual variety identification on DSL-TL dev.
CC-TLD models are trained on zero-annotation web text from mc4 filtered by
country TLD. \textbf{\#train} = number of country cells in the trained model.
English CC-TLD is omitted: the \texttt{.us} TLD is essentially unused,
making US English invisible to URL-based filtering.}
\label{tab:crosslingual}
\end{table}

Table~\ref{tab:crosslingual} reveals that the competitive margin of CC-TLD
depends critically on \textbf{training set size and variety count}.

\paragraph{Portuguese: CC-TLD wins (74.7\% vs.\ 66.0\%).}
This is the clearest success case. The mc4 corpus contains 50{,}000 Brazilian
and 14{,}647 Portuguese documents with clean TLD signal (\texttt{.br}/\texttt{.pt}),
versus only a few thousand sentences in the DSL-TL training split.
The CC-TLD model dramatically improves European Portuguese recognition
(50.1\% $\to$ 67.4\%), because European Portuguese is under-represented in
the small labeled set but well-represented on \texttt{.pt} domains.
The result mirrors our Arabic finding: \emph{scale beats label quality}.

\paragraph{Spanish: oracle wins, but CC-TLD is competitive (71.7\% vs.\ 78.4\%).}
The Spanish CC-TLD model covers \textbf{19 countries}, while the evaluation
tests only Spain (ES) vs.\ Argentina (AR).  This \emph{variety dilution} effect
hurts: texts from Spain sometimes get assigned to Mexico (\texttt{.mx}), and
Argentine texts to Chile or Colombia.  The 19-class model is penalized for
knowing too much.  A pruned 2-class model restricted to \texttt{.es}/\texttt{.ar}
would likely exceed the oracle---a direction we leave for future work.

\paragraph{English: CC-TLD is not applicable.}
American English dominates the web but almost entirely through \texttt{.com},
not \texttt{.us}.  The \texttt{.us} TLD holds a negligible fraction of US web
content, making URL-based filtering ineffective for the US/GB distinction.
The oracle model achieves 78.5\% macro-F1 on DSL-TL EN, confirming that
character $n$-grams capture the orthographic differences between American and
British English.

% ─────────────────────────────────────────────────────────────────────────────
\section{Discussion}
\label{sec:discussion}

\paragraph{Why does geography predict dialect?}
The assumption underlying our approach---that geographic origin correlates with
dialect---is well-supported empirically. \citet{eisenstein-etal-2010-latent} showed
that geographic distance predicts lexical similarity on Twitter, and the
NADI task was explicitly designed around the observation that Arabic dialects
follow country boundaries \cite{abdul-mageed-etal-2020-nadi}. Character $n$-grams
capture the phonological and orthographic fingerprints of dialectal variation
(e.g., the \textit{q/g/k} realization in Egyptian Arabic; the borrowing patterns
in Maghrebi dialects), which are consistent enough across geo-labeled data to
serve as training signal.

\paragraph{When to use geo-Twitter vs. oracle.}
Our results suggest a practical decision rule: if a labeled dialect corpus
exists but is small ($<$5{,}000 examples per variety), geo-Twitter training is
likely competitive or superior.  If labeled data is large and clean, and the
target dialect is a high-diaspora variety (Gulf states), oracle data remains
preferable.  The two sources are complementary and can be combined.

\paragraph{Limitations.}
Character $n$-gram models have a known ceiling on dialectal tasks: they capture
surface form well but not pragmatics, code-switching, or semantic meaning.
State-of-the-art systems for Arabic dialect ID use fine-tuned transformer models
(CAMeLBERT, AraBERT) that capture long-range dependencies and achieve 30--40\%
macro-F1 on NADI~\cite{abdul-mageed-etal-2020-nadi}---roughly twice our oracle
result. Geographic weak supervision could in principle be applied to transformer
pre-training corpora, which we leave for future work.

Additionally, our geo-Twitter corpus relies on self-declaration, which introduces
two forms of noise: (1) expatriates writing in a dialect different from their
declared location (especially in Gulf countries), and (2) automated or
promotional accounts that may not produce natural dialectal text.

% ─────────────────────────────────────────────────────────────────────────────
\section{Conclusion}
\label{sec:conclusion}

Geographic metadata has long been known to correlate with dialect; character
$n$-gram models have long been used for variety identification.  What this
paper provides is the first systematic answer to the question of when combining
the two is sufficient to replace dialect annotation entirely.

The answer depends on three factors we name and characterize: training data
volume, label noise from diaspora populations, and variety-set dilution.
When volume is high and noise is low---as in Portuguese CC-TLD or Arabic
geo-Twitter---geographic supervision matches or surpasses annotation at no
cost.  When label noise dominates---Gulf Arabic, where expatriates degrade the
geo-signal---annotation retains an advantage.  When evaluation is narrower
than training---Spanish evaluated on two varieties from a 19-country model---
dilution penalizes coverage breadth.

A secondary finding is the concrete documentation of the coverage failure of
fastText and GlotLID on Arabic: 14 of 21 NADI countries are invisible to
these systems, not because their dialects are hard to model, but because no
variety-level signal ever entered their training pipelines.  Geographic
filtering, at scale, repairs this gap without a single human annotation.

We release all code, trained models, and evaluation scripts.

% ─────────────────────────────────────────────────────────────────────────────
\bibliography{references}

\appendix

% ─────────────────────────────────────────────────────────────────────────────
\section{NADI 2020 Country Labels}
\label{app:labels}

The 21 countries in NADI 2020, with ISO-3166-1 alpha-2 codes and approximate
geo-Twitter training sizes: AE (25K), BH (25K), DJ (0), DZ (16K), EG (55K),
IQ (15K), JO (27K), KM (0), KW (40K), LB (27K), LY (35K), MA (11K), MR (0),
OM (18K), PS (42K), QA (30K), SA (26K), SD (14K), SO (0), SY (16K), TN (9K),
YE (10K).

\section{Hyperparameters}
\label{app:hyperparams}

All models use order-4 character $n$-gram LMs with modified Kneser-Ney smoothing.
Log-prior weight $\lambda=1.0$. Minimum training examples per country: 10.
No hyperparameter tuning was performed on the dev set.

\end{document}
